\documentclass[../comorg-part1.tex]{subfiles}

\begin{document}
\section{Instruction Set}
\begin{figure}[H]
    \centering
    \includegraphics[width=0.9\textwidth]{images/3-types-instructs.png}
    \caption*{The 3 types of instructions.}
\end{figure}

\subsection{Data Transfer Instructions}
Instructions that move data between registers and/or memory.
\subsubsection{Register Data Transfer}
\begin{itemize}
    \item MOV instruction utilises register direct or immediate addressing.
    \begin{figure}[H]
        \centering
        \includegraphics[width=0.7\textwidth]{images/mov-movs.png}
        \caption*{MOV/MOVS instruction.}
    \end{figure}
    \item MVN instruction inverts the source operand bitwise by moving
    into the destination register.
    \begin{figure}[H]
        \centering
        \includegraphics[width=0.7\textwidth]{images/movn.png}
        \caption*{MVN instruction.}
    \end{figure}
\end{itemize}

\subsubsection{Memory Data Transfer}
\begin{itemize}
    \item LDR accesses memory at effective address using indirect register
    addressing modes.
    \begin{figure}[H]
        \centering
        \includegraphics[width=0.9\textwidth]{images/ldr-instructs.png}
        \caption*{LDR/LDRB instruction.}
    \end{figure}
    \item STR copies data to effective address using indirect register addressing modes. 
    \begin{figure}[H]
        \centering
        \includegraphics[width=0.9\textwidth]{images/str-instructs.png}
        \caption*{STR/STRB instruction.}
    \end{figure}
\end{itemize}

\subsection{Data Processing Instructions}
\subsubsection{Arithmetic Instructions}
Supports register direct and immediate addressing (for rightmost operand only).
\begin{itemize}
    \item ADD/ADDS
    \begin{figure}[H]
        \centering
        \includegraphics[width=0.4\textwidth]{images/adds-instruct.png}
        \caption*{Rightmost operand can be immediate value or register.}
    \end{figure}
    \begin{figure}[H]
        \centering
        \begin{subfigure}[b]{0.45\textwidth}
            \centering
            \includegraphics[width=\textwidth]{images/twos-comp-of.png}
            \caption*{V=1: overflow when adding \textbf{signed} numbers.}
        \end{subfigure}
        \hfill
        \begin{subfigure}[b]{0.45\textwidth}
            \centering
            \includegraphics[width=\textwidth]{images/unsigned-of.png}
            \caption*{C=1: overflow when adding \textbf{unsigned} numbers.}
        \end{subfigure}
        \caption*{ADDS affects all the flags NZVC.}
    \end{figure}
    \item SUB/SUBS/RSB/RSBS
    \begin{figure}[H]
        \centering
        \includegraphics[width=0.5\textwidth]{images/subs-instruct.png}
        \caption*{SUB is not commutative; rightmost operand can be immediate value or register.}
    \end{figure}
    \begin{figure}[H]
        \centering
        \includegraphics[width=0.5\textwidth]{images/rsbs-instruct.png}
        \caption*{RSB reverses the subtraction order.}
    \end{figure}
    \begin{figure}[H]
        \centering
        \begin{subfigure}[b]{0.45\textwidth}
            \centering
            \includegraphics[width=\textwidth]{images/minuend-subtrahend.png}
            % \caption{V=1: overflow when adding \textbf{signed} numbers.}
        \end{subfigure}
        \hfill
        \begin{subfigure}[b]{0.45\textwidth}
            \centering
            \includegraphics[width=\textwidth]{images/sub-cflag.png}
            % \caption{C=1: overflow when adding \textbf{unsigned} numbers.}
        \end{subfigure}
        \caption*{SUBS affecting C flag.}
    \end{figure}
    \begin{figure}[H]
        \centering
        \begin{subfigure}[b]{0.45\textwidth}
            \centering
            \includegraphics[width=\textwidth]{images/sub-c-1.png}
            \caption{C=1 if minuend $>$ subtrahend}
        \end{subfigure}
        \hfill
        \begin{subfigure}[b]{0.45\textwidth}
            \centering
            \includegraphics[width=\textwidth]{images/sub-c-0.png}
            \caption{C=0 if minuend $<$ subtrahend.}
        \end{subfigure}
        % \caption*{SUBS affecting C flag.}
    \end{figure}
    \begin{figure}[H]
        \centering
        \includegraphics[width=0.5\textwidth]{images/sub-vflag.png}
        \caption*{V=1 when the result is out of the signed 32-bit range.}
    \end{figure}

    \item ADC/SBC/RSC [with S suffix to influence NZVC]
    \begin{figure}[H]
        \centering
        \includegraphics[width=0.8\textwidth]{images/adc-sbc-rsc.png}
        \caption*{Uses value of C flag at that point.}
    \end{figure}
\end{itemize}

\subsubsection{Logical instructions}
Boolean operators; S suffix influences N/Z flags.
\begin{itemize}
    \item MVN/MVNS
    \begin{figure}[H]
        \centering
        \includegraphics[width=0.8\textwidth]{images/mvns-instruct.png}
        \caption*{MVN takes two operands, performs bitwise inversion.}
    \end{figure}
    \item AND/ANDS
    \begin{figure}[H]
        \centering
        \includegraphics[width=0.8\textwidth]{images/and-instruct.png}
        \caption*{AND can be used to clear bits to 0.}
    \end{figure}
    \item ORR/ORRS
    \begin{figure}[H]
        \centering
        \includegraphics[width=0.8\textwidth]{images/orr-instruct.png}
        \caption*{ORR can be used to set bits to 1.}
    \end{figure}
    \item EOR/EORS
    \begin{figure}[H]
        \centering
        \includegraphics[width=0.8\textwidth]{images/eor-instruct.png}
        \caption*{EOR used to complement bits.}
    \end{figure}
    \item Logical shift left (LSL) \& Logical Shift Right (LSR)
    \begin{figure}[H]
        \centering
        \includegraphics[width=0.8\textwidth]{images/lsl-lsr.png}
        \caption*{Multiply (shift left) or divide (shift right) by factor of $2^N$, for N bits shifted.}
    \end{figure}
    LSL: For signed/unsigned multiply, '0's are shifted into the LSB of
    the register from the right. \newline
    LSR: For unsigned divide, '0's are shifted into the MSB of the register
    from the left. \newline
    'S' suffix will influence C flag.
    \begin{figure}[H]
        \centering
        \includegraphics[width=0.9\textwidth]{images/lsl-lsr-instructs.png}
        \caption*{Shift operation is applied to the rightmost operand, 
        number of bits to be shifted can be specified as immediate value or within a register.}
    \end{figure}
    \item Arithmetic shift right (ASR)
    \begin{figure}[H]
        \centering
        \includegraphics[width=0.4\textwidth]{images/asr.png}
        % \caption*{  }
    \end{figure}
    ASR: For signed divide, the sign bit is shifted into the MSB from the left.\newline
    'S' suffix will influence C flag.
    \item Rotate right (ROR) or Rotate Right Extended (RRX)
    \begin{figure}[H]
        \centering
        \begin{subfigure}[b]{0.45\textwidth}
            \centering
            \includegraphics[width=\textwidth]{images/ror.png}
            \caption*{Bits shifted out are returned in at leftmost end,
            and also placed in C flag.}
        \end{subfigure}
        \hfill
        \begin{subfigure}[b]{0.45\textwidth}
            \centering
            \includegraphics[width=\textwidth]{images/rrx.png}
            \caption*{C flag is shifted into the register at the leftmost end,
            while the bit shifted out replaces the C flag.}
        \end{subfigure}
        % \caption*{SUBS affecting C flag.}
    \end{figure}
    \begin{figure}[H]
        \centering
        \includegraphics[width=0.9\textwidth]{images/ror-instruct.png}
        \caption*{Rotate operation is applied to the rightmost operand.}
    \end{figure}
\end{itemize}

\subsection{Program Control Instructions}
To disrupt the program's normal sequential flow, the contents of the 
Program Counter (PC) is modifed, either directly or by using Branch instruction.

\medskip
\noindent If a branch is executed based on some condition(s),
it is a Conditional Branch or Branch on Conditional Code (Bcc).

\medskip
\noindent If the condition in the conditional code (cc) is true, 
a displacement is added to the PC. Note that the PC value is 
8 bytes ahead of the current Bcc being executed.

\medskip
\noindent The displacement range is $\pm32$MB.

\begin{figure}[H]
    \centering
    \includegraphics[width=0.9\textwidth]{images/bcc-table.png}
    \caption*{The 15 conditional codes. Each code will be prefixed with 'B' in the code.}
\end{figure}

Signed vs unsigned values
\begin{itemize}
    \item For testing signed values, use \textbf{GT, LT, GE, LE}.
    \item For testing unsigned values, use \textbf{HI, LO, HS, LS}.
\end{itemize}

\bigskip
\noindent Example of using different condition codes for the same function.
\begin{figure}[H]
    \centering
    \includegraphics[width=0.7\textwidth]{images/bne-loop.png}
    \caption*{BNE tests for Z=0}
\end{figure}
\begin{figure}[H]
    \centering
    \includegraphics[width=0.7\textwidth]{images/bpl-loop.png}
    \caption*{BPL test for N=0}
\end{figure}

\begin{itemize}
    \item CMP instruction subtracts source operand from destination
    register and sets the CC flags (does not 'S' suffix).
    \begin{figure}[H]
        \centering
        \begin{subfigure}[b]{0.45\textwidth}
            \centering
            \includegraphics[width=\textwidth]{images/subs-vs-cmp.png}
            \caption*{R1 is modified here.}
        \end{subfigure}
        \hfill
        \begin{subfigure}[b]{0.45\textwidth}
            \centering
            \includegraphics[width=\textwidth]{images/cmp.png}
            \caption*{R1 unchanged. CMP tests (R1 - R2).}
        \end{subfigure}
        % \caption*{SUBS affecting C flag.}
    \end{figure}
    \item CMN (Compare negative)
    \begin{figure}[H]
    \centering
    
\includegraphics[width=0.7\textwidth]{images/cmn.png}
    \caption*{Tests for (R0 + R1). Essentially ADDS without modifying the destination register.}
    \end{figure}

    \item TST (Test bits)
    \begin{figure}[H]
    \centering
    \includegraphics[width=0.7\textwidth]{images/tst.png}
    \caption*{Performs bitwise AND between the two operands, without modifying the destination register.}
    \end{figure}

    \item TEQ (Test equivalence)
    \begin{figure}[H]
    \centering
    \includegraphics[width=0.7\textwidth]{images/teq.png}
    \caption*{Performs bitwise XOR between the two operands, without modifying the destination register.}
    \end{figure}

\end{itemize}


\end{document}